\documentclass[10pt,letterpaper,final]{report}
\usepackage[margin=.75in]{geometry}
\usepackage[utf8]{inputenc}
\usepackage{amsmath}
\usepackage{amsfonts}
\usepackage{amssymb}
\usepackage{amsthm}
\renewcommand\qedsymbol{$\blacksquare$}
\usepackage{enumerate}
\usepackage{enumitem}
\usepackage{hyperref}
\usepackage{pdfpages}
\usepackage{graphics}
\usepackage{graphicx}
\usepackage{tikz}
\usepackage{tikz-qtree}
\usepackage{forest}
\usepackage{float}
\usetikzlibrary{automata,arrows}
\usetikzlibrary{shapes.multipart, positioning}
\usepackage{mdframed}
\usepackage[
  backend=biber,
  style=numeric-comp,
  sorting=none
]{biblatex}
\addbibresource{references.bib}

\usepackage{minted}
\usemintedstyle{algol_nu}
\usepackage{xltabular}
\usepackage{pdflscape}
\usepackage{tcolorbox}
\usepackage{array}
\usepackage{todonotes}
\setuptodonotes{inline}  % render \todo notes as full-width blocks in the text flow, not margin boxes

\newcommand{\reason}[1]{&& \text{#1}}

% Based on template from COSC-417 at Towson University, originally authored by Marius Zimand

% To input a script, adapt the following:
% \inputminted[breaklines, breakanywhere, fontsize=\footnotesize, frame=single, fontsize=\footnotesize]{python}{script.py}

\begin{document}
\begin{center}
  \fbox{
    \vbox{
      \begin{flushleft}
        Jonathan Llovet \\  % authors' names
        EN.625.252.81 \\  %class
        2026-10-11 \\  % date
        Module 6 - Homework 6 \\ % ID
      \end{flushleft}
      \center{\Large{\textbf{Linear Algebra}}}
      %\end{mdframed}
    } % end vbox
  } % end fbox
  \vline
\end{center}

\medskip
\noindent\textbf{Prompt 1:}
\medskip

Let $S_1$ and $S_2$ be subspaces of $R^n$. Define the union $S_1 \cup S_2$, the intersection $S_1 \cap S_2$, and the direct sum $S_1$ and $S_2$, denoted $S_1 \oplus S_2$. Of these new sets, which are and which are not subspaces of $R^n$? Give an argument or counterexample to justify your answer.

\medskip
\noindent\textbf{Answer:}
\medskip



\medskip
\noindent\textbf{Prompt 2:}
\medskip

Is the set containing just the zero vector, $S = \{\mathbf{0}\}$ is a subspace of any $R^n$? If so, find a basis and clearly argue why that basis is linearly independent. What is its dimension?

\medskip
\noindent\textbf{Answer:}
\medskip

There are some helpful sections from our textbook to help with answering this.

See \textcite[Ch. 2.9, p. 164]{LayEtAl2020LinearAlgebraAndItsApplications},

\begin{tcolorbox}[colback=blue!10!white, colframe=blue!10!black, title=Definition - Dimension of a Subspace]
  The \textbf{dimension} of a nonzero subspace $H$, denoted by $\dim H$, is the number of vectors in any basis for $H$. The dimension of the zero subspace $\{\mathbf{0}\}$ is defined to be zero.
\end{tcolorbox}

The note on this reads: ``The zero subspace has \emph{no} basis (because the zero vector by itself forms a linearly dependent set).''

There is an argument for this \textcite[Ch. 1.7, p. 61]{LayEtAl2020LinearAlgebraAndItsApplications}, "The zero vector is linearly dependent because $x_1\mathbf{0} = \mathbf{0}$ has many nontrivial solutions." In fact, there are infinitely many such solutions in $R^n, n \ge 1$.

Armed with this, we should see whether the set we are interested in can satisfy the conditions for being a subspace on \textcite[Ch. 2.8, p. 155]{LayEtAl2020LinearAlgebraAndItsApplications}

\begin{tcolorbox}[colback=blue!10!white, colframe=blue!10!black, title=Definition - Subspace]
  A \textbf{subspace} of $R^n$ is any set $H$ in $R^n$ that has three properties:
  \begin{enumerate}[label=\alph*.]
    \item The zero vector is in $H$.
    \item For each $\mathbf{u}$ and $\mathbf{v}$ in $H$, the sum $\mathbf{u} + \mathbf{v}$ is in $H$.
    \item For each $\mathbf{u}$ in $H$ and each scalar $c$, the vector $c\mathbf{u}$ is in $H$.
  \end{enumerate}
\end{tcolorbox}

By the preceding, the set $S = \{\mathbf{0}\}$ can form a subspace of any $R^n$, though by definition there is no basis, and the dimension is zero by definition. Let $H = S$.

The first condition is satisfied by the definition of $S$, i.e. $\mathbf{0} \in S$, because $S = \{\mathbf{0}\}$.

For the second condition, let $\mathbf{u} = \mathbf{0}$ and let $\mathbf{v} = \mathbf{0}$. Then $\mathbf{u} + \mathbf{v} = \mathbf{0} + \mathbf{0} = \mathbf{0} \in S$.

For the third condition, let $c$ be any scalar. Because $c0 = 0,\ \forall c$, therefore $c\mathbf{0} = \mathbf{0} \in S$.

\pagebreak
\noindent\textbf{Prompt 3:}
\medskip

List the standard basis of $R$, $R^2$, and $R^3$.  Then list two other bases for $R$, $R^2$, and $R^3$ each. Give a reason why another basis besides the standard basis might be used.  How many bases does $R$, $R^2$, and $R^3$ have?

\medskip
\noindent\textbf{Answer:}
\medskip

Standard basis vectors:

$R: \mathcal{B}_1 = \left\{\begin{bmatrix}
    1
  \end{bmatrix}\right\}$
\medskip
$R^2: \mathcal{B}_2 = \left\{\begin{bmatrix}
    1 \\
    0
  \end{bmatrix}, \begin{bmatrix}
    0 \\
    1
  \end{bmatrix}\right\}$
\medskip
$R^3: \mathcal{B}_3 = \left\{\begin{bmatrix}
    1 \\
    0 \\
    0
  \end{bmatrix}, \begin{bmatrix}
    0 \\
    1 \\
    0
  \end{bmatrix},\begin{bmatrix}
    0 \\
    0 \\
    1
  \end{bmatrix}\right\}$

Alternates:

$R: \mathcal{B}_1 = \left\{\begin{bmatrix}
    -1
  \end{bmatrix}\right\}$
\medskip
$R^2: \mathcal{B}_2 = \left\{\begin{bmatrix}
    1 \\
    2
  \end{bmatrix}, \begin{bmatrix}
    4 \\
    3
  \end{bmatrix}\right\}$
\medskip
$R^3: \mathcal{B}_3 = \left\{\begin{bmatrix}
    1 \\
    5 \\
    3
  \end{bmatrix}, \begin{bmatrix}
    2 \\
    6 \\
    1
  \end{bmatrix},\begin{bmatrix}
    1 \\
    1 \\
    1
  \end{bmatrix}\right\}$

$R: \mathcal{B}_1 = \left\{\begin{bmatrix}
    2
  \end{bmatrix}\right\}$
\medskip
$R^2: \mathcal{B}_2 = \left\{\begin{bmatrix}
    3 \\
    1
  \end{bmatrix}, \begin{bmatrix}
    1 \\
    4
  \end{bmatrix}\right\}$
\medskip
$R^3: \mathcal{B}_3 = \left\{\begin{bmatrix}
    2  \\
    10 \\
    8
  \end{bmatrix}, \begin{bmatrix}
    1 \\
    6 \\
    0
  \end{bmatrix},\begin{bmatrix}
    2 \\
    2 \\
    2
  \end{bmatrix}\right\}$

It might be reasonable to use a nonstandard set of basis vectors to have the frame of reference that you use for analysis match the way that a problem domain presents data, for example, so that elementary operations are simpler, and so that transformations are easier to represent.

For each of $R$, $R^2$, and $R^3$, there are infinitely many bases, because it is possible to find infinitely many linearly independent vectors for each.

\pagebreak

\section*{Source Code}
The full source code, including LaTeX, tooling, other artifacts, and git history is available on my Github repository for the class here:
\begin{center}
  \url{https://github.com/jllovet/jhu-en-625-252-fa26-linear-algebra}
\end{center}

\printbibliography

\end{document}
